\scanpage{67}
NB $F(\mathbb 1_{B'})$ \unsure{devient alors|devra être|devient donc} un ``objet-\deleted{($k$)}anneau'' de $B$ (``$k$-Algèbre'' quand il y a un $k$) $F$ transforme ``anneaux'' de $B'$ en ``anneau'' de $B$, en respectant les \unsure{propr.|proprs} d'unité, associativité ou commutativité.

Les $F$ en question forment une catégorie. Si on a $B''$, $B'$, $B$, ils se composent, d'où $2$-catégorie des $\otimes$-catégories $k$-additives.

\uline{Supposons maintenant que} $B'=\mathbb{K}$ -- (catégorie dans laquelle se trouvent ttes nos résolutions injectives --). Un foncteur $k$-lin.
\[
\mathbb{K}\xrightarrow{F}B
\]
s'identifie alors à un $C\in B^{\bullet}$. Donc on sait ce qu'il faut entendre par ``structure multiplicative sur le complexe de cochaînes $C^{\bullet}$'' de $B^{\bullet}$ (on dira: str. mult. demi-simpliciale'' pour ne pas prêter à confusion).
% archive sheet number: 66
